\CatchFileBetweenTags{\NcVal}{calculations/flavor_geometry.tex}{NcVal}
\CatchFileBetweenTags{\NcEq}{calculations/flavor_geometry.tex}{NcEq}
\CatchFileBetweenTags{\WeinbergGUTVal}{calculations/flavor_geometry.tex}{WeinbergGUTVal}
\CatchFileBetweenTags{\WeinbergGUTEq}{calculations/flavor_geometry.tex}{WeinbergGUTEq}
\CatchFileBetweenTags{\GammaRatioVal}{calculations/flavor_geometry.tex}{GammaRatioVal}
\CatchFileBetweenTags{\GammaRatioEq}{calculations/flavor_geometry.tex}{GammaRatioEq}
\CatchFileBetweenTags{\GammaRatioExperimentalValue}{calculations/flavor_geometry.tex}{GammaRatioExperimentalValue}
\CatchFileBetweenTags{\GammaRatioAccText}{calculations/flavor_geometry.tex}{GammaRatioAccText}
\CatchFileBetweenTags{\AreaVal}{calculations/flavor_geometry.tex}{AreaVal}
\CatchFileBetweenTags{\AreaEq}{calculations/flavor_geometry.tex}{AreaEq}

\appendix
\crefalias{section}{appendix}
\section{Geometric Audits of the Interaction Sector}

We present four structural tests demonstrating that the invariants $\mathbb{S} = \{D=4, \sigma=5, \chi=2\}$ govern not just mixing angles, but the logic of anomalies, unification, and decay widths.

\subsection{The QCD Axial Anomaly (Geometric Origin of \texorpdfstring{$N_c$}{Nc})}
In the Standard Model, the anomaly cancellation required for renormalizability mandates that the sum of charges over fermions vanishes. This works only if there are exactly $N_c=3$ colors. In the $E_8$-Persistence framework, ``Color" is identified as the \textbf{Interaction Remainder}, the degrees of freedom in the interaction geometry ($\sigma$) not consumed by the topological boundary ($\chi$).

\begin{equation}
    N_c \equiv \NcEq = \NcVal
\end{equation}

This confirms that the factor of 3 is the necessary geometric surplus required to embed a $\chi=2$ boundary within a $\sigma=5$ interaction manifold.

\subsection{The Weinberg Angle at Unification}
Standard Grand Unified Theories (like $SU(5)$) predict that at the unification scale, the Weak Mixing Angle converges to $\sin^2 \theta_W = 3/8 = 0.375$. The $E_8$ lattice recovers this naturally as the ratio of the Color Sector ($N_c=3$) to the Total Lattice Rank in the symmetric, unprojected phase. Before chiral truncation, both $D_4$ sectors are active, yielding $2D=8$ dimensions.

\begin{equation}
    \sin^2 \theta_W(M_{GUT}) = \WeinbergGUTEq = \WeinbergGUTVal
\end{equation}

This demonstrates that the low-energy value derived in Paper I ($\approx 0.223$) is the result of projecting this symmetric $3/8$ ratio onto the constrained 4D manifold.

\subsection{The Gauge Width Ratio (\texorpdfstring{$\Gamma_W / \Gamma_Z$}{GammaWDivGammaZ})}
The decay widths of the $W$ and $Z$ bosons are determined by the available geometric phase space. We define the \textbf{Aperture Ratio} as the ratio of the Interaction Symmetry ($\sigma$) to the Total Weak Aperture ($\sigma+1$).

\begin{equation}
    \mathcal{R}_{geo} = \GammaRatioEq \approx \GammaRatioVal
\end{equation}

Comparing this to the experimental ratio derived from the Particle Data Group (2024):
\begin{center}
    Experimental Ratio: \GammaRatioExperimentalValue
\end{center}

\GammaRatioAccText \ This suggests the decay width hierarchy is strictly regulated by the integer geometry of the interaction generators ($5/6$).

\subsection{The Unitarity Area (Holographic Phase)}
Standard quantum mechanics defines the Jarlskog Invariant $J$ as twice the area of the Unitarity Triangle ($J = 2A$). In the lattice model, this factor of 2 arises from the \textbf{Chiral Projection}. The area $A$ exists on the projected 4D manifold, but the invariant $J$ accounts for the full bit-depth of the lattice node ($N=32$) relative to the chiral channel ($\nu=16$).

\begin{equation}
    A_{triangle} = \AreaEq \approx \AreaVal
\end{equation}

This identifies the Unitarity Triangle as the ``half-integer" holographic projection of the CP-violating phase space.